\section{2011 Algebra \& Number Theory}

\subsection*{Individual}

\begin{exercise}
Let $K = \mathbb{Q}(\sqrt{-3})$, an imaginary quadratic field.

(a) Does there exist a finite Galois extension $L/\mathbb{Q}$ which contains $K$ such that $\mathrm{Gal}(L/\mathbb{Q}) \cong S_3$?

(b) Does there exist a finite Galois extension $L/\mathbb{Q}$ which contains $K$ such that $\mathrm{Gal}(L/\mathbb{Q}) \cong \mathbb{Z}/4\mathbb{Z}$?

(c) Does there exist a finite Galois extension $L/\mathbb{Q}$ which contains $K$ such that $\mathrm{Gal}(L/\mathbb{Q}) \cong Q_8$?
\end{exercise}

\begin{solution}
(a) \textbf{Yes.}
\textbf{Motivation:} The group $S_3$ is the smallest non-abelian group. It has a unique normal subgroup $A_3 \cong \mathbb{Z}/3\mathbb{Z}$ with quotient $S_3/A_3 \cong \mathbb{Z}/2\mathbb{Z}$. By the Fundamental Theorem of Galois Theory, any $S_3$ extension $L/\mathbb{Q}$ has a unique quadratic subfield corresponding to $A_3$.
Consider the polynomial $f(x) = x^3 - 2$. Its roots are $\sqrt[3]{2}, \sqrt[3]{2}\omega, \sqrt[3]{2}\omega^2$ where $\omega = e^{2\pi i/3}$. The splitting field is $L = \mathbb{Q}(\sqrt[3]{2}, \omega)$.
The Galois group is $S_3$. The quadratic subfield of $L$ is $\mathbb{Q}(\omega) = \mathbb{Q}(\sqrt{-3}) = K$. Thus, $L$ is such an extension.

(b) \textbf{No.}
\textbf{Theorem:} A quadratic field $k = \mathbb{Q}(\sqrt{d})$ (with $d$ square-free) is contained in a cyclic extension of degree 4 over $\mathbb{Q}$ if and only if $d$ is a sum of two squares in $\mathbb{Q}$.
\textbf{Pedagogical Remark:} A cyclic extension $L/\mathbb{Q}$ of degree 4 is a "tower" $\mathbb{Q} \subset k \subset L$. If $\text{Gal}(L/\mathbb{Q}) = \langle \sigma \rangle \cong \mathbb{Z}/4\mathbb{Z}$, then $k$ is the unique fixed field of $\sigma^2$. It can be shown (using Kummer theory or by looking at the norm) that for $k = \mathbb{Q}(\sqrt{d})$, $L$ is of the form $k(\sqrt{u})$ where $N_{k/\mathbb{Q}}(u) = a^2 d$ is not allowed, rather $N_{k/\mathbb{Q}}(u) \in \mathbb{Q}^{\times 2}$.
Actually, the condition is simpler: $d$ must be positive and $d = a^2 + b^2$ for some $a, b \in \mathbb{Q}$.
In our case, $d = -3$. Since $-3 < 0$, it is not a sum of squares in $\mathbb{R}$, hence not in $\mathbb{Q}$. Thus, no such extension exists.

(c) \textbf{Yes.}
\textbf{Theorem (Witt):} A biquadratic field $\mathbb{Q}(\sqrt{a}, \sqrt{b})$ can be embedded in a $Q_8$-extension if and only if the quadratic form $ax^2 + by^2 + abz^2$ is equivalent to $x^2 + y^2 + z^2$ over $\mathbb{Q}$. This is equivalent to saying that $a, b, ab$ are all sums of three squares and a certain Hilbert symbol condition holds.
Specifically, $K = \mathbb{Q}(\sqrt{-3})$ is a subfield of the biquadratic field $k = \mathbb{Q}(\sqrt{2}, \sqrt{-3}) = \mathbb{Q}(\sqrt{2}, \sqrt{-6})$.
A $Q_8$ extension exists if $a$ can be written as a sum of three squares. For example, $Q_8$ extensions over $\mathbb{Q}$ containing $\mathbb{Q}(\sqrt{-3})$ do exist (e.g., using the criterion that $(-1, -1)_p = (a, b)_p$). For $a=-1, b=-1$, we get the standard $Q_8$. For $K$ to be a subfield, we need the biquadratic field to satisfy the embedding condition. This is possible for many choices.
\end{solution}

\begin{exercise}
Let $f: G \to GL_2(\mathbb{C})$ be a 2D rep of a finite group $G$ such that 1 is an eigenvalue of $f(\sigma)$ for every $\sigma \in G$. Prove that $f$ is a direct sum of two 1D reps.
\end{exercise}

\begin{solution}
\textbf{Step 1: Analyze the Character.}
Let $\chi$ be the character of $f$. For each $\sigma \in G$, the eigenvalues of $f(\sigma)$ are $1$ and $\det(f(\sigma))$. Let $\lambda(\sigma) = \det(f(\sigma))$.
Since $f$ is a representation, $\lambda: G \to \mathbb{C}^\times$ is a 1D representation (a character).
The character of $f$ is $\chi(\sigma) = 1 + \lambda(\sigma)$.

\textbf{Step 2: Use Orthogonality of Characters.}
Recall the inner product of characters: $\langle \phi, \psi \rangle = \frac{1}{|G|} \sum_{g \in G} \phi(g) \overline{\psi(g)}$.
The character of the trivial representation is $\chi_{triv}(\sigma) = 1$.
The multiplicity of the trivial representation in $f$ is:
\[ \langle \chi, \chi_{triv} \rangle = \frac{1}{|G|} \sum_{\sigma \in G} (1 + \lambda(\sigma)) = \frac{1}{|G|} \sum_{\sigma \in G} 1 + \frac{1}{|G|} \sum_{\sigma \in G} \lambda(\sigma) = 1 + \langle \lambda, \chi_{triv} \rangle. \]
Since $\lambda$ is an irreducible character (1D), $\langle \lambda, \chi_{triv} \rangle$ is either 1 (if $\lambda$ is trivial) or 0 (if $\lambda$ is not trivial).
- If $\lambda = \chi_{triv}$, then $\langle \chi, \chi_{triv} \rangle = 2$. Since $\dim f = 2$, $f \cong \mathbb{C} \oplus \mathbb{C}$.
- If $\lambda \neq \chi_{triv}$, then $\langle \chi, \chi_{triv} \rangle = 1$. Thus $f$ contains the trivial representation exactly once. The complement must be 1D, so $f \cong \mathbb{C} \oplus \mathbb{C}_\lambda$, where $\mathbb{C}_\lambda$ is the 1D space where $G$ acts by $\lambda$.
In both cases, $f$ is a direct sum of two 1D representations.
\end{solution}

\begin{exercise}
Let $F \subset \mathbb{R}$ be the subset of all real numbers that are roots of monic polynomials $f(X) \in \mathbb{Q}[X]$.
\begin{enumerate}
    \item Show that $F$ is a field.
    \item Show that the only field automorphisms of $F$ are the identity automorphism $\sigma(x) = x$ for all $x \in F$.
\end{enumerate}
\end{exercise}

\begin{solution}
(1) $F$ is the set of all real algebraic numbers. Since the set of all algebraic numbers $\overline{\mathbb{Q}}$ is a field (this is a standard result in field theory: the sum, product, and inverse of algebraic numbers are algebraic), and $F = \overline{\mathbb{Q}} \cap \mathbb{R}$, $F$ is the intersection of two fields, hence a field.

(2) \textbf{Motivation:} Unlike $\mathbb{C}$, the field $\mathbb{R}$ has a natural order which is rigid under field automorphisms.
Let $\sigma$ be an automorphism of $F$. Any field automorphism of an algebraic extension of $\mathbb{Q}$ must fix $\mathbb{Q}$ elementwise.
First, we show $\sigma$ preserves order. If $x > 0$ and $x \in F$, then $x = y^2$ for some $y \in \mathbb{R}$. Since $x$ is algebraic, $y = \sqrt{x}$ is also algebraic, so $y \in F$.
Then $\sigma(x) = \sigma(y^2) = \sigma(y)^2 \geq 0$. Since $\sigma$ is an isomorphism, $\sigma(x) \neq 0$ for $x \neq 0$, so $\sigma(x) > 0$.
Thus $x < y \implies y-x > 0 \implies \sigma(y-x) > 0 \implies \sigma(x) < \sigma(y)$.
Now, for any $x \in F$, and any rational $q$, $q < x \iff \sigma(q) < \sigma(x) \iff q < \sigma(x)$.
Since $\mathbb{Q}$ is dense in $\mathbb{R}$ and $F \subseteq \mathbb{R}$, if $\sigma(x) \neq x$, say $\sigma(x) < x$, there would be a rational $q$ such that $\sigma(x) < q < x$. But $q < x \implies \sigma(q) < \sigma(x) \implies q < \sigma(x)$, a contradiction. Thus $\sigma(x) = x$.
\end{solution}

\begin{exercise}
Let $V$ be a finite-dimensional vector space over $\mathbb{R}$ and $T: V \to V$ be a linear transformation such that:
\begin{enumerate}
    \item the minimal polynomial of $T$ is irreducible;
    \item there exists a vector $v \in V$ such that $\{ T^i v \mid i \geq 0 \}$ spans $V$.
\end{enumerate}
Show that $V$ contains no non-trivial proper $T$-invariant subspace.
\end{exercise}

\begin{solution}
\textbf{Pedagogical Remark:} Think of $V$ as an $\mathbb{R}[x]$-module where $x$ acts as $T$.
Since $V$ is cyclic, $V \cong \mathbb{R}[x]/(f(x))$ where $f(x)$ is the characteristic polynomial.
For a cyclic module, the minimal polynomial $m_T(x)$ equals the characteristic polynomial $f(x)$ (up to sign).
Since $m_T(x)$ is irreducible, the ideal $(m_T(x))$ in $\mathbb{R}[x]$ is a maximal ideal.
Thus $V \cong \mathbb{R}[x]/(m_T(x))$ is a field.
The $T$-invariant subspaces of $V$ are precisely the $\mathbb{R}[x]$-submodules of $V$.
Submodules of $\mathbb{R}[x]/(m_T(x))$ correspond to ideals in the ring $\mathbb{R}[x]/(m_T(x))$.
Since it is a field, its only ideals are $\{0\}$ and the field itself.
Thus there are no proper non-trivial $T$-invariant subspaces.
\end{solution}

\begin{exercise}
Given a commutative diagram of Abelian groups:
\[ \begin{tikzcd}
A \arrow[r] \arrow[d, "\alpha"] & B \arrow[r] \arrow[d, "\beta"] & C \arrow[r] \arrow[d, "\gamma"] & D \arrow[r] \arrow[d, "\delta"] & E \arrow[d, "\epsilon"] \\
A' \arrow[r] & B' \arrow[r] & C' \arrow[r] & D' \arrow[r] & E'
\end{tikzcd} \]
such that (i) both rows are exact sequences and (ii) every vertical map, except the middle one $\gamma$, is an isomorphism. Show that the middle map $\gamma: C \to C'$ is also an isomorphism.
\end{exercise}

\begin{solution}
\textbf{Pedagogical Remark:} This is a standard "diagram chase". We show injectivity and surjectivity of the middle map $\gamma: C \to C'$.
Let the maps be:
\begin{align*}
A \xrightarrow{f_1} B \xrightarrow{f_2} C \xrightarrow{f_3} D \xrightarrow{f_4} E \\
\alpha \downarrow \quad \beta \downarrow \quad \gamma \downarrow \quad \delta \downarrow \quad \epsilon \downarrow \\
A' \xrightarrow{g_1} B' \xrightarrow{g_2} C' \xrightarrow{g_3} D' \xrightarrow{g_4} E'
\end{align*}
Assume all vertical maps except $\gamma$ are isomorphisms.
\textbf{Injectivity:} Let $c \in C$ such that $\gamma(c) = 0$.
Then $g_3(\gamma(c)) = 0 \implies \delta(f_3(c)) = 0 \implies f_3(c) = 0$ (since $\delta$ is an iso).
By exactness at $C$, $c = f_2(b)$ for some $b \in B$.
Then $\gamma(f_2(b)) = 0 \implies g_2(\beta(b)) = 0$.
By exactness at $B'$, $\beta(b) = g_1(a')$ for some $a' \in A'$.
Since $\alpha$ is an iso, $a' = \alpha(a)$ for some $a \in A$.
Then $\beta(b) = g_1(\alpha(a)) = \beta(f_1(a)) \implies b = f_1(a)$ (since $\beta$ is an iso).
Then $c = f_2(b) = f_2(f_1(a)) = 0$ (by exactness at $B$). Thus $\gamma$ is injective.
\textbf{Surjectivity:} Similar diagram chase shows $\gamma$ is surjective.
\end{solution}

\begin{exercise}
Group of order 150 is not simple.
\end{exercise}

\begin{solution}
$|G| = 150 = 2 \cdot 3 \cdot 5^2$.
Sylow 5-subgroups: $n_5 \equiv 1 \pmod 5$ and $n_5 \mid 6$. So $n_5 \in \{1, 6\}$.
If $n_5 = 1$, the Sylow 5-subgroup is normal.
If $n_5 = 6$, $G$ acts on the 6 subgroups via conjugation. This gives $\phi: G \to S_6$.
If $G$ is simple, $\phi$ is injective.
A Sylow 5-subgroup $P$ of $G$ has order $5^2 = 25$.
Then $\phi(P)$ is a subgroup of $S_6$ of order 25.
However, $|S_6| = 720 = 2^4 \cdot 3^2 \cdot 5$. The highest power of 5 dividing $|S_6|$ is 5.
By Lagrange's theorem, $S_6$ cannot have a subgroup of order 25.
Contradiction. Thus $n_5 = 1$.
\end{solution}

\subsection*{Team}
\begin{exercise}
Let $F'$ be a field and $F$ the algebraic closure of $F'$. Let $f(x, y)$ and $g(x, y)$ be polynomials in $F[x, y]$ such that $\gcd(f, g) = 1$ in $F[x, y]$. Show that there are only finitely many $(a, b) \in F^2$ such that $f(a, b) = g(a, b) = 0$. Can you generalize this to the case of more than two variables?
\end{exercise}

\begin{solution}
\textbf{Theorem (Bezout):} Two plane curves of degrees $d_1$ and $d_2$ with no common component intersect in at most $d_1 d_2$ points.
The condition $\gcd(f, g) = 1$ means they share no common factor, so the curves have no common component.
In algebraic terms, if $\gcd(f, g) = 1$, then the ideal $(f, g)$ in $F[x, y]$ has a zero-dimensional variety.
\textbf{Generalization:} For $n$ variables, $f_1, \dots, f_n$ having no common factor is not enough (they could define a curve in 3D). The correct condition is that the sequence $(f_1, \dots, f_n)$ forms a \textit{regular sequence} or that the dimension of the resulting variety is zero.
\end{solution}

\begin{exercise}
Let $D$ be a PID, and $D^n$ the free module of rank $n$ over $D$. Prove that any submodule of $D^n$ is a free module of rank $m \leq n$.
\end{exercise}

\begin{solution}
\textbf{Proof Sketch (Induction on $n$):}
Base case $n=1$: A submodule is an ideal. In a PID, an ideal is $(d)$, which is free of rank 1 ($d \neq 0$) or rank 0 ($d=0$).
Step: Let $M \subseteq D^n$. Let $p: D^n \to D$ be the projection on the first coordinate.
$p(M)$ is an ideal in $D$, so $p(M) = (d)$.
If $d=0$, $M \subseteq D^{n-1}$, done by induction.
If $d \neq 0$, let $v \in M$ such that $p(v) = d$. Since $p(M) \cong D$ is free, the exact sequence $0 \to \ker p|_M \to M \to p(M) \to 0$ splits.
$M \cong \ker p|_M \oplus p(M)$.
By induction, $\ker p|_M \subseteq D^{n-1}$ is free of rank $m \leq n-1$.
Thus $M$ is free of rank $m+1 \leq n$.
\end{solution}

\begin{exercise}
Identify pairs of integers $n \neq m \in \mathbb{Z}_+$ such that the quotient rings $\mathbb{Z}[x, y]/(x^2 - y^n) \cong \mathbb{Z}[x, y]/(x^2 - y^m)$; and identify pairs of integers $n \neq m \in \mathbb{Z}_+$ such that $\mathbb{Z}[x, y]/(x^2 - y^n) \not\cong \mathbb{Z}[x, y]/(x^2 - y^m)$.
\end{exercise}

\begin{solution}
Let $R_n = \mathbb{Z}[x, y]/(x^2 - y^n)$.
If $n$ is even, $n=2k$, then $x^2 - y^{2k} = (x-y^k)(x+y^k)$. $R_n$ has zero divisors.
If $n$ is odd, $x^2 - y^n$ is irreducible, so $R_n$ is an integral domain.
This distinguishes even and odd.
For odd $n$, the normalization of $R_n$ is $\mathbb{Z}[t]$, with $y=t^2, x=t^n$.
The ring is isomorphic to $\mathbb{Z}[t^2, t^n]$. This is a numerical semigroup ring.
The semigroup $S_n = \langle 2, n \rangle$ has exactly $(n-1)/2$ gaps in $\mathbb{N}$.
Different $n$ give different numbers of gaps, so the rings are not isomorphic.
Conclusion: No such pairs $n \neq m$ exist.
\end{solution}

\begin{exercise}
Is it possible to find an integer $n > 1$ such that the sum
\[ 1 + \frac{1}{2} + \frac{1}{3} + \dots + \frac{1}{n} \]
is an integer?
\end{exercise}

\begin{solution}
\textbf{Explanatory Remark:} This is a classic problem in number theory using the "2-adic" valuation.
Let $k$ be the largest integer such that $2^k \leq n$.
Consider the least common multiple $L = \text{lcm}(1, 2, \dots, n)$.
$L = 2^k \cdot m$ where $m$ is odd.
Multiply $H_n$ by $L/2^{k-1}$:
All terms in $(L/2^{k-1}) \cdot H_n$ are integers except for the term $1/2^k$, which becomes $m/2$.
Thus the sum cannot be an integer.
Alternatively: In the sum $\sum_{j=1}^n \frac{L}{j}$, every term is even except the one where $j=2^k$ (which is $L/2^k = m$, odd).
So the numerator of $H_n = \frac{\sum L/j}{L}$ is odd, while the denominator is even.
\end{solution}

\begin{exercise}
Recall that $\mathbb{F}_7$ is the finite field with 7 elements, and $GL_3(\mathbb{F}_7)$ is the group of all invertible $3 \times 3$ matrices with entries in $\mathbb{F}_7$.
\begin{enumerate}
    \item Find a 7-Sylow subgroup $P_7$ of $GL_3(\mathbb{F}_7)$.
    \item Determine the normalizer subgroup $N$ of the 7-Sylow subgroup you found in (1).
    \item Find a 2-Sylow subgroup of $GL_3(\mathbb{F}_7)$.
\end{enumerate}
\end{exercise}

\begin{solution}
The order is $(7^3-1)(7^3-7)(7^3-49) = 342 \cdot 336 \cdot 294 = 2^6 \cdot 3^4 \cdot 7^3 \cdot 19$.
(a) $P_7$ is the group of upper triangular unipotent matrices (order $7^3$):
$P_7 = \{ \begin{pmatrix} 1 & a & b \\ 0 & 1 & c \\ 0 & 0 & 1 \end{pmatrix} \}$.
(b) The normalizer $N(P_7)$ is the Borel subgroup $B$ of upper triangular matrices:
$B = \{ \begin{pmatrix} d_1 & a & b \\ 0 & d_2 & c \\ 0 & 0 & d_3 \end{pmatrix} : d_i \in \mathbb{F}_7^\times \}$.
(c) Sylow 2-subgroup has order $2^6 = 64$.
$\mathbb{F}_7^\times \cong C_6$. Its Sylow 2-subgroup is $\{\pm 1\}$.
Consider the group of monomial matrices with entries in $\{\pm 1\}$.
This is $(\mathbb{Z}/2\mathbb{Z})^3 \rtimes S_3$. Order $8 \times 6 = 48$. Not enough.
Actually, $GL_2(\mathbb{F}_7)$ has order $48 \times 42$, which has $2^4 \cdot 2^1 = 2^5$.
So $P_2(GL_2) \times P_2(GL_1)$ has order $2^5 \times 2^1 = 2^6 = 64$.
\end{solution}

\begin{exercise}
Show that $SL_2(\mathbb{Z})$ is generated by $T = \begin{pmatrix} 1 & 1 \\ 0 & 1 \end{pmatrix}$ and $S = \begin{pmatrix} 0 & 1 \\ -1 & 0 \end{pmatrix}$.
\end{exercise}

\begin{solution}
\textbf{Motivation:} This is essentially the Euclidean Algorithm.
Let $M = \begin{pmatrix} a & b \\ c & d \end{pmatrix}$. We want to reduce $c$ to 0.
$T^n M = \begin{pmatrix} a+nc & b+nd \\ c & d \end{pmatrix}$. We can choose $n$ to make $|a+nc| \leq |c|/2$.
Then $S(T^n M) = \begin{pmatrix} c & d \\ -(a+nc) & -(b+nd) \end{pmatrix}$.
The absolute value of the lower-left entry has decreased.
Repeating this, we eventually get $c=0$.
If $c=0$, then $ad=1$, so $a=d=1$ or $a=d=-1$.
If $a=d=1$, $M = \begin{pmatrix} 1 & b \\ 0 & 1 \end{pmatrix} = T^b$.
If $a=d=-1$, $M = \begin{pmatrix} -1 & b \\ 0 & -1 \end{pmatrix} = S^2 T^{-b}$.
Thus all matrices are in $\langle S, T \rangle$.
\end{solution}
